% section: 極形式と複素平面

\section{極形式と複素平面}\label{節：極形式と複素平面}
本節では,座標平面における回転,移動,縮小,拡大およびド・モアブルの定理,オイラーの公式について扱う.

  \subsection{複素平面と極形式}
    $i^2=-1$を満たす虚数単位$i$を用いると,実数$x,y$に対して複素数は
    \[
      z=x+iy
    \]
    と表される.
    この複素数$z$を座標平面上の点$(x,y)$に対応させたものを複素平面という.
    横軸を実軸,縦軸を虚軸とする.

    点$(x,y)$と原点との距離は
    \[
      \rho=\sqrt{x^2+y^2}
    \]
    である.
    また,正の$x$軸から点$(x,y)$へ向かう半直線まで反時計回りに測った角を$\theta$とすると,
    \[
      x=\rho\cos\theta,\qquad
      y=\rho\sin\theta
    \]
    と表せる.
    したがって,\,$z$は
    \[
      z=\rho(\cos\theta+i\sin\theta)
    \]
    と表される.
    この表し方を複素数の極形式という.

    $\rho$は$z$の絶対値であり,\,$|z|=\rho$と書く.
    $z\ne0$のとき,\,$\theta$を$z$の偏角という.
    偏角には$2\pi$の整数倍を加えても同じ点を表すため,一意には定まらない.
    $z=0$では$\rho=0$となり,偏角は定まらない.

    例えば,\,$\rho=2,\,\theta=\dfrac{\pi}{3}$の点は
    \[
      (x,y)
      =
      \left(
        2\cos\frac{\pi}{3},
        2\sin\frac{\pi}{3}
      \right)
      =
      (1,\sqrt{3})
    \]
    である.
    直交座標では$(1,\sqrt{3})$,極形式では
    \[
      2\left(
        \cos\frac{\pi}{3}
        +i\sin\frac{\pi}{3}
      \right)
    \]
    と表される.

  \subsection{複素数による平面上の変換}
    複素数の加法と乗法は,座標平面上の点の移動や回転,拡大・縮小に対応する.

    まず,\,$b=u+iv$を固定した複素数とし,\,$z=x+iy$に$b$を加える.
    \[
      w=z+b=(x+u)+i(y+v)
    \]
    となるので,点$(x,y)$は点$(x+u,y+v)$へ移る.
    これは,すべての点を同じベクトル$(u,v)$だけ平行移動する操作である.

    次に,\,$z=\rho(\cos\theta+i\sin\theta)$に
    \[
      c=s(\cos\phi+i\sin\phi)
      \qquad(s>0)
    \]
    を掛ける.
    三角関数の加法定理より,
    \begin{align*}
      cz
        &=s\rho
          (\cos\phi+i\sin\phi)
          (\cos\theta+i\sin\theta)\\
        &=s\rho
          \{\cos(\theta+\phi)+i\sin(\theta+\phi)\}.
    \end{align*}
    したがって,絶対値は$s$倍され,偏角は$\phi$だけ増える.

    特に,\,$s=1$なら原点を中心とする角$\phi$の回転である.
    $s>1$なら原点からの距離が$s$倍される拡大,\,$0<s<1$なら原点からの距離が$s$倍される縮小となる.
    例えば,\,$i=\cos\dfrac{\pi}{2}+i\sin\dfrac{\pi}{2}$であるから,\,$z$に$i$を掛けると原点を中心に反時計回りに$\dfrac{\pi}{2}$だけ回転する.

    回転・拡大縮小の中心を$z_0$とし,移動先の中心を$w_0$とすると,これらをまとめた変換は
    \[
      w=c(z-z_0)+w_0
    \]
    と表せる.
    まず中心からの差$z-z_0$に$c$を掛けて回転・拡大縮小し,その後$w_0$を加えて中心を移す.

  \subsection{ド・モアブルの定理}
    まず,極形式どうしの積を考える.
    三角関数の加法定理から,
    \begin{align*}
      &(\cos\alpha+i\sin\alpha)
       (\cos\beta+i\sin\beta)\\
      &\quad=
       \cos\alpha\cos\beta-\sin\alpha\sin\beta\\
      &\qquad
       +i(\sin\alpha\cos\beta+\cos\alpha\sin\beta)\\
      &\quad=
       \cos(\alpha+\beta)+i\sin(\alpha+\beta)
    \end{align*}
    となる.
    つまり,絶対値が$1$の複素数どうしを掛けると,偏角が加わる.

    これを繰り返すと,任意の整数$n$について
    \[
      (\cos\theta+i\sin\theta)^n
      =
      \cos(n\theta)+i\sin(n\theta)
    \]
    が成り立つ.
    これをド・モアブルの定理という.

    正の整数$n$については,数学的帰納法で確かめられる.
    $n=1$では明らかであり,\,$n$で成り立つと仮定すると,
    \begin{align*}
      (\cos\theta+i\sin\theta)^{n+1}
        &=
        \{\cos(n\theta)+i\sin(n\theta)\}
        (\cos\theta+i\sin\theta)\\
        &=
        \cos((n+1)\theta)+i\sin((n+1)\theta)
    \end{align*}
    となる.
    また,\,$n=0$では両辺とも$1$である.
    $n<0$の場合は,加法定理より
    \[
      (\cos\theta+i\sin\theta)
      (\cos(-\theta)+i\sin(-\theta))=1
    \]
    であるから,
    \[
      (\cos\theta+i\sin\theta)^{-1}
      =
      \cos(-\theta)+i\sin(-\theta)
    \]
    を用いて正の整数の場合に帰着できる.

    例えば,
    \[
      \left(
        \cos\frac{\pi}{6}
        +i\sin\frac{\pi}{6}
      \right)^6
      =
      \cos\pi+i\sin\pi
      =
      -1
    \]
    である.
    このように極形式で表すと,累乗では絶対値を累乗し,偏角を指数倍すればよい.

  \subsection{オイラーの公式}
    次に,複素数の指数表示$e^{i\theta}$を極限を用いて考える.
    ここでは
    \[
      e^{i\theta}
      =
      \lim_{n\to\infty}
      \left(1+\frac{i\theta}{n}\right)^n
    \]
    と定める.
    前節で準備した極限を用いて,この値を求める.

    $\theta\ne0$とし,\,$\phi_n$を
    \[
      -\frac{\pi}{2}<\phi_n<\frac{\pi}{2},
      \qquad
      \tan\phi_n=\frac{\theta}{n}
    \]
    を満たす角とする.
    また,
    \[
      r_n=\sqrt{1+\frac{\theta^2}{n^2}}
    \]
    とおくと,
    \[
      \cos\phi_n=\frac{1}{r_n},
      \qquad
      \sin\phi_n=\frac{\theta}{nr_n}
    \]
    である.
    よって,
    \[
      1+\frac{i\theta}{n}
      =
      r_n(\cos\phi_n+i\sin\phi_n)
    \]
    と極形式で表せる.

    まず,\,$r_n^n$の極限を求める.
    自然対数の基本極限より,
    \begin{align*}
      \log(r_n^n)
        &=
        \frac{n}{2}
        \log\left(1+\frac{\theta^2}{n^2}\right)\\
        &=
        \frac{\theta^2}{2n}
        \frac{
          \log\left(1+\frac{\theta^2}{n^2}\right)
        }{
          \frac{\theta^2}{n^2}
        }\\
        &\longrightarrow0.
    \end{align*}
    したがって,
    \[
      r_n^n\longrightarrow1
    \]
    である.

    次に,\,$n\phi_n$の極限を求める.
    $\tan\phi_n=\dfrac{\theta}{n}\to0$より$\phi_n\to0$であり,
    \[
      n\phi_n
      =
      \theta\frac{\phi_n}{\tan\phi_n}
      \longrightarrow\theta
    \]
    となる.
    ここでは,\,$\lim_{x\to0}\dfrac{x}{\tan x}=1$を用いた.

    ド・モアブルの定理により,
    \[
      \left(1+\frac{i\theta}{n}\right)^n
      =
      r_n^n
      \{\cos(n\phi_n)+i\sin(n\phi_n)\}
    \]
    である.
    $r_n^n\to1,\,n\phi_n\to\theta$と三角関数の連続性から,
    \[
      e^{i\theta}
      =
      \lim_{n\to\infty}
      \left(1+\frac{i\theta}{n}\right)^n
      =
      \cos\theta+i\sin\theta
    \]
    を得る.
    $\theta=0$の場合も両辺が$1$であり,同じ式が成り立つ.
    これをオイラーの公式という.

    オイラーの公式により,複素数の極形式
    \[
      z=\rho(\cos\theta+i\sin\theta)
    \]
    は
    \[
      z=\rho e^{i\theta}
    \]
    とも表せる.
    さらに,ド・モアブルの定理は
    \[
      z^n
      =
      \rho^n e^{in\theta}
      =
      \rho^n\{\cos(n\theta)+i\sin(n\theta)\}
      \qquad(n\in\mathbb Z)
    \]
    と書ける.

  \subsection{複素数の実数乗}
    $z=\rho e^{i\theta}$で$\rho>0$とし,偏角$\theta$を一つ選んで固定する.
    この選び方に基づいて,実数$x$に対する複素数のべき乗を
    \[
      z^x=\rho^x e^{ix\theta}
      =
      \rho^x\{\cos(x\theta)+i\sin(x\theta)\}
    \]
    と定めることができる.

    ただし,偏角は$\theta+2\pi k$とも表せるため,整数でない$x$に対する$z^x$の値は,どの偏角を選ぶかによって変わる.
    $x$が整数なら,選ぶ偏角によらず同じ値になる.
    また,\,$x=\dfrac{m}{n}$を有理数とし,\,$n\ge1$として上の定義で得た値を$n$乗すると,
    \[
      (z^{m/n})^n=z^m
    \]
    となる.
    したがって,この定義で得られる$z^{m/n}$は$z^m$の$n$乗根の一つである.

    \subsection{議論のまとめ}
    ここまでの議論が長くなったので,扱った内容をまとめる.

    複素数$z=x+iy$を座標平面上の点$(x,y)$に対応させると,複素数を平面上の点として考えられる.
    極形式
    \[
      z=\rho(\cos\theta+i\sin\theta)
    \]
    では,\,$\rho$が原点からの距離を表し,\,$\theta$が偏角を表す.

    複素数に一定の数を加えることは平行移動に対応し,極形式で表した複素数を掛けることは回転と拡大縮小に対応する.
    積を繰り返すと偏角が加わるため,ド・モアブルの定理が得られる.
    また,第\ref{節：極限}節で準備した極限を用いるとオイラーの公式
    \[
      e^{i\theta}=\cos\theta+i\sin\theta
    \]
    が導かれる.
    このように,複素数を極形式で表すことで,座標平面上の変換と複素数の計算を結び付けられる.